% ============================================================================
\chapter{开发环境搭建}
\label{ch:environment}
% ============================================================================

\section{本章目标}

在动手写任何内核代码之前，先把工具链准备好：交叉编译器、汇编器、
链接器、GRUB 打包工具和 QEMU 模拟器。

\section{宿主机环境}

\begin{itemize}
  \item 操作系统：Ubuntu 22.04+ / WSL2（本书使用 WSL2 + VS Code Remote）
  \item 编辑器：VS Code + clangd 扩展
\end{itemize}

\section{必要工具安装}

% TODO: 从 docs/howto.md 整理
\begin{lstlisting}[style=shellstyle]
sudo apt update
sudo apt install -y nasm gcc binutils grub-pc-bin \
    xorriso qemu-system-x86 gcc-multilib gdb bear
\end{lstlisting}

各工具的作用：

\begin{longtable}{ll}
  \toprule
  \textbf{工具} & \textbf{用途} \\
  \midrule
  \texttt{nasm}           & x86 汇编器（NASM 语法） \\
  \texttt{gcc -m32}       & C 编译器（32 位模式） \\
  \texttt{ld}             & 链接器 \\
  \texttt{grub-mkrescue}  & 生成 GRUB 可引导 ISO \\
  \texttt{xorriso}        & ISO 9660 工具（grub-mkrescue 依赖） \\
  \texttt{qemu-system-i386} & x86 模拟器 \\
  \texttt{gdb}            & 调试器 \\
  \texttt{bear}           & 生成 \texttt{compile\_commands.json}（给 clangd） \\
  \bottomrule
\end{longtable}

\section{项目目录结构}

% TODO: 从 copilot-instructions.md 的文件结构复制
\begin{lstlisting}[style=shellstyle]
my_microkernel/
  src/
    boot/          # Multiboot2 引导 + 链接脚本
    include/       # 公共头文件
    kernel/        # 内核实现
  docs/            # 项目文档
  book/            # 本书 LaTeX 源码
  Makefile         # 顶层 wrapper
\end{lstlisting}

\section{构建与运行闭环}

% TODO: 详细展开 make iso / make run / make debug 流程

\begin{lstlisting}[style=shellstyle]
cd src
make iso        # 编译内核 + 打包 ISO
make run        # QEMU 启动
make debug      # QEMU -S 等待 GDB 连接
make gdb        # 另一个终端启动 GDB
\end{lstlisting}

\section{VS Code 调试集成}

% TODO: 从 docs/debug.md 整理 launch.json / tasks.json 配置

\section{本章小结}

工具链就绪后，下一章我们写第一行代码——Multiboot2 引导头和内核入口。
